% GENERATED FILE. Edit canonical lessons, then run npm run generate:books.
% canonical-source-hash: fnv1a64:8b965c7c53486d5c
% canonical-lessons: RU-C500-vyazat, RU-C500-spitsy, RU-C500-klubok, RU-C500-vyshivat, RU-C500-kruzhevo, RU-R500-review-493-500

\chapter{Knitting and Embroidery}
\label{ch:ru-knitting-and-embroidery}
\hlchaptermodality{500}

\begin{chapteropening}
\textbf{By the end of this chapter:} \emph{I can talk about knitting, embroidery and lace.}

Complete the last lesson of chapter 500: I can talk about knitting, embroidery and lace.
\end{chapteropening}

\section[\texorpdfstring{\ru{вязать} (vyazát)}{вязать (vyazát)}]{\ru{вязать} (vyazát) — to knit; to tie}

\label{lesson:RU-C500-vyazat}



\begin{quote}
Before the new one: say the Russian for a tea party, then the Russian for to stay as a guest.
\end{quote}

\subsection*{You'll want to know: \ru{вязать}}
\textbf{\ru{вязать}} (\emph{vyazát}) — \textquotedblleft{}to knit; to tie\textquotedblright{}.

\begin{quote}
\ru{Бабушка} \ru{вяжет} \ru{мне} \ru{тёплый} \ru{шарф}.
\end{quote}

\begin{quote}
\emph{Bábushka vyázhet mne tyóplyy sharf.}
\end{quote}

\begin{quote}
Grandma is knitting me a warm scarf.
\end{quote}

Imperfective; present \ru{вяжу}, \ru{вяжешь} (\ru{з} becomes \ru{ж}). Perfective \ru{связать}.

\begin{cousinweb}
\emph{\ru{вязать}} is inherited from Proto-Slavic \emph{vęzati} 'to tie, to bind' (EtymDB).
\end{cousinweb}

\begin{grammarlens}[title={What you've built}]
One more word for telling, explaining and describing at B1.
\end{grammarlens}

\subsection*{Your turn}
\begin{itemize}
\raggedright
  \item \emph{Say it:} \emph{vyazát}
  \item \emph{Say it:} \emph{vyazát}, once more
  \item \emph{Say it:} a sentence of your own with \emph{vyazát}
  \item \emph{Recall:} say the Russian for a tea party, then the Russian for to stay as a guest, then say \emph{vyazát} again
\end{itemize}

\begin{tcolorbox}[breakable,colback=teal!4,colframe=teal!35!black,title={Before you move on}]
What does \ru{вязать} mean? (\textquotedblleft{}To knit; to tie\textquotedblright{}.) Say it once more.
\end{tcolorbox}
\section[\texorpdfstring{\ru{спицы} (spítsy)}{спицы (spítsy)}]{\ru{спицы} (spítsy) — knitting needles}

\label{lesson:RU-C500-spitsy}



\begin{quote}
Before the new one: say the Russian for to stay as a guest, then the Russian for to knit.
\end{quote}

\subsection*{You'll want to know: \ru{спицы}}
\textbf{\ru{спицы}} (\emph{spítsy}) — \textquotedblleft{}knitting needles\textquotedblright{}.

\begin{quote}
\ru{Для} \ru{этого} \ru{свитера} \ru{нужны} \ru{толстые} \ru{спицы}.
\end{quote}

\begin{quote}
\emph{Dlya étogo svítera nuzhný tólstyye spítsy.}
\end{quote}

\begin{quote}
You need thick needles for this sweater.
\end{quote}

Plural of \emph{\ru{спица}}, which also means a spoke of a wheel.

\begin{cousinweb}
\emph{\ru{спицы}} is the plural of \emph{\ru{спица}} 'a spoke, a knitting needle', which EtymDB traces to Proto-Slavic \emph{st\ru{ъ}pica}; EtymDB has no entry for the plural.
\end{cousinweb}

\begin{grammarlens}[title={What you've built}]
One more word for telling, explaining and describing at B1.
\end{grammarlens}

\subsection*{Your turn}
\begin{itemize}
\raggedright
  \item \emph{Say it:} \emph{spítsy}
  \item \emph{Say it:} \emph{spítsy}, once more
  \item \emph{Say it:} a sentence of your own with \emph{spítsy}
  \item \emph{Recall:} say the Russian for to stay as a guest, then the Russian for to knit, then say \emph{spítsy} again
\end{itemize}

\begin{tcolorbox}[breakable,colback=teal!4,colframe=teal!35!black,title={Before you move on}]
What does \ru{спицы} mean? (\textquotedblleft{}Knitting needles\textquotedblright{}.) Say it once more.
\end{tcolorbox}
\section[\texorpdfstring{\ru{клубок} (klubók)}{клубок (klubók)}]{\ru{клубок} (klubók) — a ball of yarn}

\label{lesson:RU-C500-klubok}



\begin{quote}
Before the new one: say the Russian for to knit, then the Russian for knitting needles.
\end{quote}

\subsection*{You'll want to know: \ru{клубок}}
\textbf{\ru{клубок}} (\emph{klubók}) — \textquotedblleft{}a ball of yarn\textquotedblright{}.

\begin{quote}
\ru{Котёнок} \ru{играет} \ru{с} \ru{клубком} \ru{шерсти}.
\end{quote}

\begin{quote}
\emph{Kotyónok igráyet s klubkóm shérsti.}
\end{quote}

\begin{quote}
The kitten is playing with a ball of wool.
\end{quote}

Masculine; genitive \ru{клубка} (the vowel drops).

\begin{cousinweb}
\emph{\ru{клубок}} is built on \emph{\ru{клуб}} 'a ball, a puff (of smoke)' with \emph{-\ru{ок}}; EtymDB lists \emph{\ru{клубок}} without a source.
\end{cousinweb}

\begin{grammarlens}[title={What you've built}]
One more word for telling, explaining and describing at B1.
\end{grammarlens}

\subsection*{Your turn}
\begin{itemize}
\raggedright
  \item \emph{Say it:} \emph{klubók}
  \item \emph{Say it:} \emph{klubók}, once more
  \item \emph{Say it:} a sentence of your own with \emph{klubók}
  \item \emph{Recall:} say the Russian for to knit, then the Russian for knitting needles, then say \emph{klubók} again
\end{itemize}

\begin{tcolorbox}[breakable,colback=teal!4,colframe=teal!35!black,title={Before you move on}]
What does \ru{клубок} mean? (\textquotedblleft{}A ball of yarn\textquotedblright{}.) Say it once more.
\end{tcolorbox}
\section[\texorpdfstring{\ru{вышивать} (vyshivát)}{вышивать (vyshivát)}]{\ru{вышивать} (vyshivát) — to embroider}

\label{lesson:RU-C500-vyshivat}



\begin{quote}
Before the new one: say the Russian for knitting needles, then the Russian for a ball of yarn.
\end{quote}

\subsection*{You'll want to know: \ru{вышивать}}
\textbf{\ru{вышивать}} (\emph{vyshivát}) — \textquotedblleft{}to embroider\textquotedblright{}.

\begin{quote}
\ru{Мама} \ru{вышивает} \ru{цветы} \ru{на} \ru{подушке}.
\end{quote}

\begin{quote}
\emph{Máma vyshiváyet tsvetý na podúshke.}
\end{quote}

\begin{quote}
Mum is embroidering flowers on a cushion.
\end{quote}

Imperfective; perfective \ru{вышить} (\ru{вышью}). The noun is \emph{\ru{вышивка}} — embroidery.

\begin{cousinweb}
\emph{\ru{вышивать}} is built on \emph{\ru{вы}-} and \emph{\ru{шить}} 'to sew'; EtymDB lists the verb without a source.
\end{cousinweb}

\begin{grammarlens}[title={What you've built}]
One more word for telling, explaining and describing at B1.
\end{grammarlens}

\subsection*{Your turn}
\begin{itemize}
\raggedright
  \item \emph{Say it:} \emph{vyshivát}
  \item \emph{Say it:} \emph{vyshivát}, once more
  \item \emph{Say it:} a sentence of your own with \emph{vyshivát}
  \item \emph{Recall:} say the Russian for knitting needles, then the Russian for a ball of yarn, then say \emph{vyshivát} again
\end{itemize}

\begin{tcolorbox}[breakable,colback=teal!4,colframe=teal!35!black,title={Before you move on}]
What does \ru{вышивать} mean? (\textquotedblleft{}To embroider\textquotedblright{}.) Say it once more.
\end{tcolorbox}
\section[\texorpdfstring{\ru{кружево} (krúzhevo)}{кружево (krúzhevo)}]{\ru{кружево} (krúzhevo) — lace}

\label{lesson:RU-C500-kruzhevo}



\begin{quote}
Before the new one: say the Russian for a ball of yarn, then the Russian for to embroider.
\end{quote}

\subsection*{You'll want to know: \ru{кружево}}
\textbf{\ru{кружево}} (\emph{krúzhevo}) — \textquotedblleft{}lace\textquotedblright{}.

\begin{quote}
\ru{На} \ru{ней} \ru{было} \ru{белое} \ru{платье} \ru{с} \ru{кружевами}.
\end{quote}

\begin{quote}
\emph{Na ney býlo béloye plátye s kruzhevámi.}
\end{quote}

\begin{quote}
She was wearing a white dress with lace.
\end{quote}

Neuter; usually plural \emph{\ru{кружева}}.

\begin{cousinweb}
EtymDB lists \emph{\ru{кружево}} without a source; it is usually traced to the root of \emph{\ru{круг}} 'circle', after the round patterns of lace.
\end{cousinweb}

\begin{grammarlens}[title={What you've built}]
One more word for telling, explaining and describing at B1.
\end{grammarlens}

\subsection*{Your turn}
\begin{itemize}
\raggedright
  \item \emph{Say it:} \emph{krúzhevo}
  \item \emph{Say it:} \emph{krúzhevo}, once more
  \item \emph{Say it:} a sentence of your own with \emph{krúzhevo}
  \item \emph{Recall:} say the Russian for a ball of yarn, then the Russian for to embroider, then say \emph{krúzhevo} again
\end{itemize}

\begin{tcolorbox}[breakable,colback=teal!4,colframe=teal!35!black,title={Before you move on}]
What does \ru{кружево} mean? (\textquotedblleft{}Lace\textquotedblright{}.) Say it once more.
\end{tcolorbox}
\section[(dialogue)]{Words from chapters 493–500 — a review}

\label{lesson:RU-R500-review-493-500}



\begin{quote}
Say the Russian for in love and a kiss.
\end{quote}

\subsection*{Your turn}
Say these aloud:

\begin{itemize}
\raggedright
  \item in love, a kiss, to flirt, tenderness, affectionate
  \item a spouse, marriage, infidelity, to be unfaithful, a parting
  \item pregnancy, a maternity hospital, maternity leave, a pram, to crawl
  \item a nappy, a dummy, a lullaby, to lull to sleep, to swaddle
  \item upbringing, to spoil, to obey, to be naughty, disobedient
  \item to be friends, friendship, a pal, someone of the same age, a classmate
  \item to visit, hospitable, hospitality, a tea party, to stay as a guest
  \item to knit, knitting needles, a ball of yarn, to embroider, lace
\end{itemize}

\begin{tcolorbox}[breakable,colback=teal!4,colframe=teal!35!black,title={Before you move on}]
In love? (\textbf{\ru{влюблённый}}.) A kiss? (\textbf{\ru{поцелуй}}.) To flirt? (\textbf{\ru{флиртовать}}.) Tenderness? (\textbf{\ru{нежность}}.) Affectionate? (\textbf{\ru{ласковый}}.) A spouse? (\textbf{\ru{супруг}}.) Marriage? (\textbf{\ru{брак}}.) Infidelity? (\textbf{\ru{измена}}.) To be unfaithful? (\textbf{\ru{изменять}}.) A parting? (\textbf{\ru{расставание}}.) Pregnancy? (\textbf{\ru{беременность}}.) A maternity hospital? (\textbf{\ru{роддом}}.) Maternity leave? (\textbf{\ru{декрет}}.) A pram? (\textbf{\ru{коляска}}.) To crawl? (\textbf{\ru{ползать}}.) A nappy? (\textbf{\ru{подгузник}}.) A dummy? (\textbf{\ru{соска}}.) A lullaby? (\textbf{\ru{колыбельная}}.) To lull to sleep? (\textbf{\ru{убаюкивать}}.) To swaddle? (\textbf{\ru{пеленать}}.) Upbringing? (\textbf{\ru{воспитание}}.) To spoil? (\textbf{\ru{баловать}}.) To obey? (\textbf{\ru{слушаться}}.) To be naughty? (\textbf{\ru{шалить}}.) Disobedient? (\textbf{\ru{непослушный}}.) To be friends? (\textbf{\ru{дружить}}.) Friendship? (\textbf{\ru{дружба}}.) A pal? (\textbf{\ru{приятель}}.) Someone of the same age? (\textbf{\ru{ровесник}}.) A classmate? (\textbf{\ru{одноклассник}}.) To visit? (\textbf{\ru{навещать}}.) Hospitable? (\textbf{\ru{гостеприимный}}.) Hospitality? (\textbf{\ru{гостеприимство}}.) A tea party? (\textbf{\ru{чаепитие}}.) To stay as a guest? (\textbf{\ru{гостить}}.) To knit? (\textbf{\ru{вязать}}.) Knitting needles? (\textbf{\ru{спицы}}.) A ball of yarn? (\textbf{\ru{клубок}}.) To embroider? (\textbf{\ru{вышивать}}.) Lace? (\textbf{\ru{кружево}}.)
\end{tcolorbox}
